\chapter{认知领域与主体边界}

\section{领域所有者}

顶层认知领域按照稳定身份、修订、生命周期、查询语义和写入权限划分。领域名称描述语义所有权，不要求一一对应独立进程、代码包或数据库。

\begin{ReportWideTable}{>{\raggedright\arraybackslash}p{3.0cm}>{\raggedright\arraybackslash}p{10.4cm}}
\captionof{table}{顶层认知领域}\label{tab:domains}
\toprule
\textbf{领域} & \textbf{职责}\\
\midrule
主体基础 & 主体标识、证据基础、跨域引用、生产者标识和共同时间/版本语义。\\
记忆 & 长期记忆、经历片段、日记、认知图式、标签和关联证据。\\
自我 & 主体对自身身份、角色、能力、倾向、价值和叙事连续性的长期认知。\\
社会认知 & 主体视角下的人际与群体关系、社会角色以及具有社会作用域的语言惯例。\\
动机 & 愿望、目标、关切和面向未来状态的优先关系。\\
认知运行时 & 会话、活动认知、工作上下文、注意分配、即时评价、驻留集和调用方工作集。\\
\bottomrule
\end{ReportWideTable}

经历维护、日记维护和离线认知属于跨领域工作流。它们可以读取多个领域并产生提案，但不形成独立的超级权威。

\section{领域对象与领域边界}

一个领域可以拥有多种一等对象。对象是否需要独立身份，由其生命周期、修订、写入和查询合同决定，而不是由术语数量决定。记忆域因此可以同时拥有记忆、经历片段、日记和认知图式；社会认知可以同时拥有关系认知和语言惯例。

领域间通过稳定引用连接。引用必须指向对象和必要的修订身份，不通过复制另一领域的完整状态建立一致性。来源发生修订、撤回、抑制或清除时，依赖对象根据自身合同进入重新核验、冲突或失效状态。

\section{主体供给与能力配置}

主体创建时由宿主或独立客户端分配 SubjectId，并按所需能力组合供给对应认知领域。未启用的领域没有隐式默认状态，也不通过空表、空插件或占位对象模拟存在。

主体初始化可以包含认知种子，用于建立初始自我认知和倾向。系统与行为宪制由宿主或 Heptalogos 的规范权威维护，不通过认知种子取得行为上的最高权限。

\section{主体分支}

主体分支创建新的 SubjectId。分支可以共享分支点之前的不可变来源和认知历史，但分支后的对象头、对象纪元、运行时状态、服务投影和管理生命周期相互独立。

恢复同一主体的备份属于身份连续性的恢复操作；克隆或分支属于创建新主体。二者在管理接口和审计记录中必须区分。
